\uuid{C24q}
\exo7id{7492}
\titre{exo7 7492}
\auteur{mourougane}
\organisation{exo7}
\datecreate{2021-08-10}
\difficulte{}
\chapitre{Géométrie affine dans le plan et dans l'espace}
\cpge{mpsi-mp2i.sous-espaces-affines, ptsi.espaces-vectoriels.b, mpsi-mp2i.matrices-applications-lineaires.a, pcsi.matrices-applications-lineaires.a, ptsi.matrices-applications-lineaires.a, mpsi-mp2i.matrices-applications-lineaires.c, pcsi.matrices-applications-lineaires.d, ptsi.matrices-applications-lineaires.d, bcpst1.applications-lineaires-matrices, bcpst2.applications-lineaires-matrices.a}
\module{Géométrie}
\niveau{L2}

\contenu{
\question{
Dans l'espace affine $E$ réel, muni d'un repère cartésien $(O, \vec{\imath},\vec{\jmath},\vec{k})$, soit $f$ une application affine dont la partie linéaire a pour matrice 
$$\begin{pmatrix}-1&0&0\cr
0&2&3\cr 0&2&3\end{pmatrix}.$$
Que peut-on dire de l'ensemble des points fixes de $f$ ?
}
}
